% sec5c_certified_drl.tex - Certified DRL null-results study (DD-033)
\subsection{Frozen-Set Evaluation of Direct DRL Phase Policies}
\label{sec:results-drl}

Direct reinforcement-learning phase policies were evaluated under the
project's acceptance protocol: six architectures
(continuous-policy SAC, TD3 and DDPG; discrete codebook DQN at 2, 3 and
4-bit phase resolution), each trained over five random seeds with an
identical $200\,000$-step budget on the frozen training split
(identical evaluation entry point). The validation-best seed of each agent was subsequently
evaluated on the frozen $500$-channel test set through the shared
\texttt{evaluate.py} entry point (\S\ref{sec:results-ao}); all rates
report the same ZF sum-rate convention as every other row.

\paragraph{Null result.}
No architecture achieves a positive mean gain over the frozen random
baseline ($53.18$~bps/Hz). The test-set deltas range from
$-0.10$ (TD3, DQN@2bit) to $+0.04$~bps/Hz (SAC), every 95\% bootstrap
interval covers zero, and every paired Wilcoxon test yields
$p \geq 0.37$
(Table~\ref{tab:certified-drl}); training-time evaluation plateaus
over all seeds corroborate the plateau at random-level values
(per-agent training-band diagnostics \texttt{fig\_band\_*} accompany
the artifact set). We therefore \emph{fail to find} evidence
that model-free reinforcement learning discovers the RIS phase
structure of \eqref{eq:sum_rate} under episodic $T{=}1$ bandit
feedback, and we report this as a null result rather than tuning
toward it.

\paragraph{Protocol qualifications.} This null is certified for the
stated protocol, and stated as such: (i)~the agents observe a
partial-CSI stack---$|h_{r}|$, $\angle h_{r}$, $|G|$, $|h_{d}|$---in
which $\angle G$ and $\angle h_{d}$ are absent, so the null applies to
policies under this observation design; a full-CSI observation mode is
implemented and released for follow-up campaigns
(supplementary decision log).
% DD-044: committed quantification of the observation-contract scope,
% artifact results/baselines/closed_form_full_csi.json (external-review round 2).
To size this qualification, we evaluate a no-learning closed-form alignment
policy with \emph{full} phase access on the same frozen channels: it reaches
$56.944 \pm 1.884$~bps/Hz ($+3.77$ over random) while the restricted-observation
DRL checkpoints are statistically at the random floor (rate gap
$\delta \in [-0.10, +0.04]$~bps/Hz across thirty committed
checkpoints)---so the null
certifies the observation contract, not gradient-based phase control per se.
% DD-044b: full-observation retraining campaign (pre-submission, per external
% review round 2). 30/30 runs completed; artifacts ckpt/full_*.pt +
% results/full/{agent}/seedS/eval_rates.json; independent cross-machine
% re-score parity 0.0000 (results/full/verification_summary.json).
The null was then tested against its own strongest objection: all thirty
checkpoints were retrained under a full-observation stack---$|h_r|$,
$\angle h_r$, $|G|$, $\angle G$, $|h_d|$, $\angle h_d$---handing every agent
the phases the DD-004 contract hides. The result is identical: the best
full-CSI family reaches $53.204 \pm 0.037$ ($+0.03$ over random,
$t \approx 0.3$), with every other family within $\delta \in [-0.11, -0.01]$
of the floor. Training computes, exploration computes, and the numbers do
not move. The certified null is therefore conditional on nothing: direct
model-free RL fails on this one-step bandit even with complete phase
information, while unfolding---which inlines the same channel into six
gradient steps---closes $97.4\%$ of the AO reference. The contrast is
between the optimization structure each method can express, not between
the information each method can see.
(ii)~In-training monitoring used the frozen evaluation set for
visualization; the loss and the selection rule that produced the
trained policies never consumed it for optimization, and the unfolding
model (whose deployment numbers headline this work) was retrained
under an audited train/validation protocol with the test set evaluated
exactly once (supplementary decision log).
(iii)~Committed per-channel evaluation artifacts cover all thirty runs
(six architectures $\times$ five seeds), as required by the acceptance
protocol.
(iv)~Disclosure: during external review of the first submission
version, an erroneous complex conjugation in the ZF pairing of the
shared rate function was identified (the SINR was evaluated as
$|\langle h, w^{*}\rangle|^{2}$ rather than
$|\langle h, w\rangle|^{2}$, a global phase rotation that had no effect
on the optimizer set but shifted every absolute rate). All baselines,
the AO reference, the unfolding retrain family, the DRL re-scoring, and
every derived table in this paper were regenerated under the corrected
pairing; in addition, the DRL evaluation wrapper was hardened to the
deterministic-action mode required by the fixed-order evaluation
contract. The reported null is a re-certification, not the original
pre-fix numbers, which are retained in the archive with a
\texttt{\_zfbuggy} suffix.

\begin{table}[t]
\centering\small
\caption{Frozen-test-set evaluation of all six DRL agents
(validation-best seed per agent; paired statistics vs.\ the frozen
random baseline over the same 500 channels).}
\label{tab:certified-drl}
\begin{tabular}{lrrrrr}
\toprule
Agent & Mean & $\Delta$ vs.\ random & $p$ & 95\% CI & \%ch.$>0$ \\
\midrule
SAC      & 53.22 & $+0.04$ & $0.72$ & $[-0.20,\,+0.27]$ & 51 \\
TD3      & 53.08 & $-0.10$ & $0.49$ & $[-0.34,\,+0.13]$ & 51 \\
DDPG     & 53.08 & $-0.09$ & $0.52$ & $[-0.33,\,+0.14]$ & 50 \\
DQN@2bit & 53.08 & $-0.10$ & $0.37$ & $[-0.33,\,+0.13]$ & 48 \\
DQN@3bit & 53.12 & $-0.06$ & $0.55$ & $[-0.28,\,+0.17]$ & 48 \\
DQN@4bit & 53.19 & $+0.02$ & $0.83$ & $[-0.20,\,+0.23]$ & 48 \\
\bottomrule
\end{tabular}
\end{table}

\paragraph{The contrast that makes the null informative.}
Exploitable head demonstrably exists in this environment: selecting
\emph{two hundred random} phase vectors per channel and keeping the best
yields $57.14$~bps/Hz ($+3.97$, $p=1.7 \times 10^{-83}$), deep
unfolding reaches $62.53$ ($+9.35$, $p=1.3 \times 10^{-83}$,
$2.6\%$ short of the Riemannian AO reference $64.19$), while every trained
policy remains statistically indistinguishable from a single random
draw. The gap between what a 200-draw random search deftly extracts and
what 30 gradient-informed training runs could not extract is the
central empirical observation of this study, and it supports the
model-driven thesis: the search landscape of phase optimization is
exploited by structure-informed optimization (unfolding, AO), not by
model-free exploration within a $T{=}1$ bandit feedback loop.

\paragraph{Why null results belong here.}
All numbers in this subsection are emitted from committed raw JSON
artifacts by the project's figure pipeline and recorded as a numbered
entry in the supplementary decision log; no post-hoc selection, seed replacement, or metric
re-weighting was performed after certification. A diagnosis of the
failure modes (reward decomposition, observation-contract coverage,
exploration coverage vs.\ the best-of-200 oracle) is the subject of a
companion technical note deferred out of the results flow; it does not
alter the reported outcome stated in
Table~\ref{tab:certified-drl}.
